\section{Native Rust inside a conjunctive-query language}
\label{sec:host-embedding}
The classical CQ account has an unusually strong separation of concerns:
an atom constrains a valuation, and the conjunction of those constraints
does not specify how to enumerate satisfying valuations. MiniLinq preserves
that account for its relational fragment. An embedded Rust expression adds
another language, with lexical binding, a type system, control flow, and
observable evaluation. Its presence changes the justification of planning,
even when the surface clause still resembles a relational selection.

Our methodological starting point is a \emph{multi-language semantics}:
make the boundary between the query language and the host explicit and
state what crosses it. Matthews and Findler formulate language composition
using such boundaries, separating properties of each language from the
properties of their interaction~\cite{host-multilanguage}. Here the boundary
passes a represented valuation into Rust and receives a value, a Boolean,
or pattern bindings. It also admits the host observations of
Section~\ref{sec:scope}. The boundary below is semantic notation for existing
Rust leaves, not an additional IR or a proposed expression evaluator.

\subsection{Six objects that must not be conflated}
\label{sec:hostobjects}
\begin{center}
\small
\begin{tabular}{p{0.20\linewidth}p{0.34\linewidth}p{0.36\linewidth}}
\toprule
Object & Question answered & What it does not establish \\
\midrule
Token tree & Which tokens and delimiters reached the macro? & Which
identifiers denote values or bind other identifiers. \\
Rust syntax tree & Which expressions, patterns, paths, and scopes were
written? & Rust name resolution, typing, or effects. \\
Resolved binding & Which declaration does an occurrence denote? & Whether
evaluating its use is pure, total, or duplicable. \\
Logical requirements & Which query variables must be available for this
clause? & How often, or on which candidate valuations, its Rust runs. \\
Evaluation domain & Which dynamic clause occurrences arise, and what
earlier filtering permits them? & Whether two occurrences commute. \\
Movement premise & Which administrative exchanges are licensed? & A
general effect analysis of arbitrary Rust. \\
\bottomrule
\end{tabular}
\end{center}

The implemented parser crosses only the first boundary. CQ constructors
use derived parsing; their host leaves use \code{syn::Expr}, \code{syn::Pat},
and \code{syn::Type}. Syn may parse \code{f(x)}, \code{|x| x + y},
\code{match} expressions, and \code{m!(tokens)} as Rust expressions without
making them alternative CQ clause forms. In particular, a macro invocation
contains tokens whose grammar is chosen by that macro; those tokens are not
automatically recursively parsed Rust expressions~\cite{host-synmacro}.

Thus \emph{arbitrary host syntax} means delegating expression syntax to the
host grammar, rather than recognizing a finite list of convenient functions
or operators. It does not mean every parsed expression is well typed in
every generated context. For example, \code{await}, \code{return}, and
borrowing expressions impose requirements on their enclosing Rust context.
Nor does parseability establish a successful CQ lowering: negation and
aggregation parse into existing CQ constructors but presently have no
implemented semantic lowering rules.

\subsection{The semantic boundary carries more than a free-variable set}
Let $\Gamma_Q$ assign types and representation modes to logical variables,
and let $\Gamma_H$ describe host bindings visible at the eventual expansion
site. Write
\[
 \Gamma_H;\Gamma_Q\vdash_H e:\tau
 \qquad\text{and}\qquad
 \rho\approx_{\Gamma_Q}\nu
\]
for host well-typedness in the elaborated context and agreement of a Rust
environment with a logical valuation. The second judgment incorporates the
value relation of Section~\ref{sec:scope}: an atom field is borrowed, whereas
a computed binding can be an owned Rust value. It is stronger than saying
that the same identifier strings are present. With the host evaluation
relation, these judgments justify an application of the boundary:
\[
\Rule{H-Eval}
 {\Gamma_H;\Gamma_Q\vdash_H e:\tau\quad
  \rho\approx_{\Gamma_Q}\nu\quad
  \langle e,\rho,\sigma\rangle\Downarrow_H(v,\sigma',\alpha)}
 {\langle\mathsf{host}[e],\nu,\rho,\sigma\rangle
  \Downarrow(v,\sigma',\alpha)}.
\]
Panic and divergence cross this boundary too. The display defines their
successful counterpart, not a rule that erases unsuccessful observations.
The implementation delegates the first judgment to rustc; its syntactic
contract checks do not derive it.

A relational expression interpretation is a further specialization. Under
$\mathcal H_{\rm rel}$, the boundary induces a total function
$\widehat e$ on the relevant valuations, independent of store and concrete
representation. Then a guard has the familiar meaning
$\sigma_{\widehat e}(L)$, and a binding extends each valuation with its
computed value. Outside that specialization, neither an arbitrary function
call nor an overloaded comparison is justified by the set-theoretic
selection laws alone. Rust explicitly permits expression effects and
trait-defined operator behavior~\cite{host-rustexpr}.

\subsection{Lexical scope, resolution, and conservative liveness}
The helper in \code{cq/rust.rs} operates before Rust name resolution.
Its result $(F,o)$ consists of visible free \emph{simple value-name
occurrences} and an opacity bit. Its free-name set is not a set of resolved
declaration identities. For example, \code{accept(x)} contributes both
\code{accept} and \code{x}; only intersection with the logical environment
makes \code{x} a query input when \code{accept} is an external function.
Conversely, a logical variable holding a callable value can make
\code{accept} itself a query input. A field name in \code{x.field}, a method
name in \code{x.method()}, and components of \code{module::function} are
not simple value-variable occurrences.

Binding-sensitive traversal matters even for a single leaf:
\begin{lstlisting}
{ let x = 1; (|witness: i32| witness + x)(*witness) }
\end{lstlisting}
The free-name set is $\{\mathit{witness}\}$. The block-local \code{x}
shadows any logical \code{x}. The closure's parameter shadows
\code{witness} only inside its body; the application argument
\code{*witness} still refers to the outer binding. Collecting identifier
tokens, or subtracting every pattern name from the whole expression, gives
the wrong answer. The FV rules in Section~\ref{sec:languages} instead
thread lexical scopes: an initializer and a let-else branch precede the
new local binding; match-arm bindings scope their guard and body;
\code{for} bindings scope only the loop body; let-chain bindings extend
later conditions and the consequent, but not the else branch. These
distinctions follow Rust's binding scopes~\cite{host-rustscope}.

For a logical environment $B$ and all syntactically logical names $Q$, the
two consumer operations are intentionally different:
\[
\begin{aligned}
 \operatorname{Live}(e,B)&=
   \begin{cases}B&o,\\F\cap B&\neg o,\end{cases}&
 \operatorname{Known}(e,Q)&=
   \begin{cases}\emptyset&o,\\F\cap Q&\neg o.\end{cases}
\end{aligned}
\]
The first is the implementation's \code{dependencies}; the second is
\code{known_uses}. An opaque macro conservatively keeps every \emph{already
available} query variable. Empty Known means the early diagnostic has no
sound syntactic claim to make, not that the macro is closed. Retaining $B$
does not invent a missing future binder or prove that a macro's identifiers
resolve as intended. Attributes, local items, and verbatim syntax also set
opacity. Local items require this care because their scope and capture
behavior differ from sequential local bindings.

Even transparent syntax leaves a resolution obligation. Syntactic
identifier patterns can denote constants in contexts decided by Rust, and
an unqualified function spelling can be shadowed. The lexical upper-bound
lemma therefore assumes the intended binding resolution. Treating that
assumption as a theorem of Syn would be an unsound strengthening.

\subsection{Logical matching and directional computation}
The following source fragment exposes the boundary in a join plan. Assume
the declared columns have type \code{i32}, \code{accept} is a host function,
and its macro counterpart is intentionally opaque.
\begin{lstlisting}
answer(x, z) :- R(x, witness),
    let y: i32 = { let x = 1; *witness + x },
    S(y, z), if accept(z), if accept!(z);
\end{lstlisting}
The dependency analysis derives these exact facts. Provided means a
positive atom's available logical columns, including equality keys; it
does not mean that all those names are fresh.
\begin{center}
\begin{tabular}{llll}
\toprule
Clause & Required & Provided & Movable \\
\midrule
$R(x,\mathit{witness})$ & $\emptyset$ & $\{x,\mathit{witness}\}$ & yes \\
$\mathsf{let}\ y=\cdots$ & $\{\mathit{witness}\}$ & $\{y\}$ & no \\
$S(y,z)$ & $\{y\}$ & $\{y,z\}$ & yes \\
$\mathsf{if}\ \mathit{accept}(z)$ & $\{z\}$ & $\emptyset$ & no \\
$\mathsf{if}\ \mathit{accept}!(z)$ & $\{x,\mathit{witness},y,z\}$ & $\emptyset$ & no \\
\bottomrule
\end{tabular}
\end{center}
The required $y$ on $S$ is essential. A value produced by an earlier
\code{let} must be computed before the atom matches it. In contrast,
$R(x,y),S(y,z)$ has no distinguished producer of $y$: either atom can
supply a partial valuation, and their common variable is an equality
constraint. Converting every shared variable into a directed
first-occurrence dependency would incorrectly forbid classical join
reordering. Removing the let-origin dependency would allow an atom to
replace a computation with an unconstrained value.

These distinctions extend inside atoms. The syntactic convention
\code{variable} recognizes an unparenthesized single-identifier path as a
logical-variable position. A new bare name binds; an already-bound or
repeated name constrains equality; \code{_} has no named binding. A general
expression is a computed matching value, not an additional producer of its
free variables. When $k$ is newly introduced by the atom,
\code{S(k, next(k))} supplies $k$ locally and has no external requirement
for it. If an earlier let binds $k$, that binding is required instead.
For a logical $k$ already available in the prefix, \code{S(next(k), z)}
requires that prefix binding; an ambient Rust name alone is not a logical
dependency. A literal position has the existing fresh-binding normalization
$R(7,x)\rightsquigarrow\mathsf{let}\ k{:}\tau=7, R(k,x)$.
That equation cannot be generalized to arbitrary computations without an
evaluation-placement premise.

Formally the dependency summary is the existing
$\Sigma;B;T\vdash c_i\triangleright(R_i,O_i,m_i,a_i,K_i)$ of
Section~\ref{sec:planning}. Atom hyperedges remain sets of shared logical
variables. Requirements add directional availability conditions, and
barriers add ordering conditions. These are different relations on the
same source occurrences; the code stores their summaries in
\code{ClauseFacts}, not in a new graph language. Atom producers of a shared
name are alternatives, so availability is more precise as the condition
$R_i\subseteq B$ than as one arbitrarily chosen producer edge per name.

\subsection{Availability does not establish movement or multiplicity}
\label{sec:hostmovement}
For pending clauses $U$, the current C-Ready test checks availability and
preserves the order of every barrier:
\[
 R_i\subseteq B,
 \qquad
 \forall j\in U.\ j<i\Longrightarrow(m_i\land m_j).
\]
Consequently two clauses may exchange written positions only inside a
region whose clauses all satisfy D-Move. This rule is implemented by
\code{ClauseState::require_ready}; a strategy selects among its candidates.
The finite D-Move recognizer admits simple atoms, literal bindings, and
recognized comparisons whose operands are syntactically primitive
\code{i32} at every occurrence. It is an explicitly delimited optimization
domain, not inference of arbitrary Rust effects.

There is a second dependence beyond the value names in $R_i$: the
\emph{domain and multiplicity of evaluation}. Suppose $R$ contains
$(1,10)$ and $(1,20)$, and the body is
\begin{lstlisting}
R(x, witness), let y: i32 = tick(x)
\end{lstlisting}
The Rust expression uses only $x$, but its prefix has two distinct
valuations. Projecting away witness first collapses them under set semantics
and changes two calls into one. Likewise a guard can suppress calls that
would have occurred before it. Ordinary liveness protects the values needed
by a computation; preservation of its evaluation domain additionally
protects witnesses and control dependence. This explains the prefix
algorithm's retained witness columns, including anonymous atom positions.
It does not by itself establish equivalence of different materialization
boundaries.

To state the missing semantic premise precisely, let $\mathcal T_c$ be a
clause's operational transformer on an enumerated valuation stream and
store, with observations including demand, host events, and termination.
An exchange requires
\[
 \mathcal T_d\circ\mathcal T_c
 \ \approx_{\rm obs}\ 
 \mathcal T_c\circ\mathcal T_d
\]
on the reachable states where both orders are well scoped. Availability is
one premise of that equation. An effect analysis can provide sufficient
conditions concerning store conflicts; classical effect systems use
conservative effect information to constrain scheduling~\cite{host-effects}.
But even nonconflicting store accesses do not automatically justify changing
the number of calls, suppressing a possible panic, or moving work before
the first cursor demand. Our equation is a semantic requirement for a
proposed optimization, not a claim that the present compiler checks it.

\paragraph{Syntactic argument order is not runtime argument order.}
The A-* argument judgments walk syntax left to right to determine scope.
Their generated operations have more than one evaluation placement:
A-Prefix computes a matching key before the atom's join; A-Local installs
a selection after that join. For
\begin{lstlisting}
answer(x) :- R(x, f(x), g());
\end{lstlisting}
$f(x)$ needs this atom's new $x$ and takes A-Local, while $g()$ takes
A-Prefix. The generated plan therefore calls $g$ on the initial singleton
prefix before it calls $f$ on matching candidates. With rows
$(1,1,100),(2,2,100)$, $f(x)=x$, and $g()=100$, full evaluation calls $g$
once and $f$ twice; the two $f$ events may occur in either input enumeration
order. The executable \code{argument_placement} example in
\code{doc/formal/example} confirms these facts for both backends, checking
the partial order without requiring a particular order of the two $f$ calls.
Preserving the order of A-Local selections does not make these
placements into Rust's left-to-right function-argument evaluation. An atom
is a DSL matching construct, not a Rust function application. A general
source effect contract must specify its intended matching-event schedule
before either placement can be judged a correct implementation of it.

\subsection{Residual code, quotation, and binding preservation}
MiniLinq inspects the relational structure while carrying most Rust leaves
forward as residual code. At lower levels those leaves occur inside
generated storage operations and iterator closures. This is a staging
boundary: construction and analysis of code occur during macro expansion;
execution of residual expressions occurs in the compiled program, at the
build or polling boundary selected by the backend. An identifier's
presence in emitted tokens does not make its value available during macro
expansion.

The language-integrated-query literature makes the host/quoted-language
distinction explicit. Cheney, Lindley, and Wadler use quotation and
normalization to give a compilability theorem for a specified query
calculus~\cite{host-linq}; Cheney et al. relate effect-based and
quotation-based query calculi by translations~\cite{host-effective}.
Those results motivate explicit boundary judgments here. They do not
imply that arbitrary native Rust normalizes to relational algebra, or that
our token-based macros implement typed open quotation. Existing
quotation and continuation mechanisms are specified separately in
Section~\ref{sec:staging}.

Finally, code movement must preserve \emph{binding}, not only spelling.
Hygienic macro systems formalize binding with syntax contexts rather than
bare strings; Flatt's scope-set account is one precise example
~\cite{host-hygiene}. Rust procedural macros have their own, explicitly
unhygienic interface~\cite{macros}. Our helper's normalized strings
serve logical-name comparison, and fresh generated names avoid known
collisions, but neither fact proves preservation under every Rust macro
expansion, shadowing context, or closure capture. The required compiler
lemma is that each residual occurrence resolves to its intended logical or
host binding and retains its typing/ownership interpretation after every
context change. This is an open host-elaboration obligation, separate from
the proved lexical traversal lemma and the finite relational theorems.
